
When \(A\) is a field, the equivalence of \textup{(i)} and
\textup{(iii)} is the definition of nilpotence, while the equivalence
of \textup{(ii)} and \textup{(iii)} is the familiar consequence of
Engel's theorem \textup{(Bourbaki, \emph{Groupes et algèbres de Lie},
Chapter I, \S4, No.~2)}. In the general case,
\textup{(iii)} plainly implies \textup{(ii)}, and \textup{(ii)}
implies \textup{(i)} by the preceding result and the fact that
\textup{(ii)} is preserved under passage to quotients and under
localization. It remains to prove that \textup{(i)} implies
\textup{(iii)}.

Suppose first that \(A\) is local Artinian, with maximal ideal
\(\mathfrak m\). Choose \(n>0\) such that \(\mathfrak m^n=0\), and
choose \(N\) such that condition \textup{(iii)} holds for
\[
  \mathfrak d(s)=\mathfrak d\otimes_A A/\mathfrak m.
\]
Taking \(N'=nN\), one sees at once that \(N'\) satisfies
\textup{(iii)} for \(\mathfrak d\).

If \(A\) is Noetherian, there are finitely many points
\(s_i\in\operatorname{Spec}(A)\) and ideals of definition
\(\mathfrak q_i\subset A_{s_i}\) such that the natural map
\[
  \mathfrak d\longrightarrow
  \prod_i\bigl(\mathfrak d\otimes_A A_i\bigr),
  \qquad A_i=A_{s_i}/\mathfrak q_i,
\]
is injective. By the preceding case, for each \(i\) there is an
integer \(N_i\) satisfying \textup{(iii)} for the Lie algebra
\(\mathfrak d\otimes_A A_i\). Taking \(N\) to be the largest of the
\(N_i\), condition \textup{(iii)} follows for \(\mathfrak d\).
Finally, the general case reduces to the Noetherian case by the
usual procedure described in EGA~IV~8.

\paragraph{\textbf{Definition 2.2.}}
\SGARefTarget{sga3:xiv:2.2}
Let \(S\) be a prescheme and let \(\mathfrak d\) be a
quasi-coherent Lie algebra over \(S\) whose underlying module is of
finite presentation. We say that \(\mathfrak d\) is
\emph{locally nilpotent} if, for every \(s\in S\), the Lie algebra
\(\mathfrak d(s)\) over \(k(s)\) is nilpotent. We say that
\(\mathfrak d\) is \emph{strictly locally nilpotent} if it is
locally free and, on every open subset \(U\) of \(S\) over which it
has constant rank \(r\), its Killing polynomial reduces to
\[
  P_{\mathfrak d}(t)=t^r.
\]

Of course, when \(\mathfrak d\) is a Lie algebra over \(S\) that is
locally free as a module, its Killing polynomial is defined as a
polynomial
\[
  P_{\mathfrak d}\in A[t],
  \qquad
  A=\Gamma\!\left(\operatorname{Sym}_{\mathcal O_S}(\mathfrak d)\right)
   \simeq \Gamma\!\left(W(\mathfrak d)\right),
\]
where \(A\) is the ring of sections of the structure sheaf of the
vector bundle \(W(\mathfrak d)\) defined by \(\mathfrak d\).

The two notions introduced in \SGARefLink{sga3:xiv:2.2}{2.2} are plainly stable under base
change and local for the fpqc topology. We record:

\paragraph{\textbf{Proposition 2.3.}}
\SGARefTarget{sga3:xiv:2.3}
\emph{If \(\mathfrak d\) is strictly locally nilpotent, then it is
locally nilpotent. The converse holds when \(\mathfrak d\) is
locally free and \(S\) is reduced.}

The proof is immediate.

\paragraph{\textbf{Definition 2.4.}}
\SGARefTarget{sga3:xiv:2.4}
Let \(S\) be a prescheme, let \(\mathfrak g\) be a Lie algebra over
\(S\) whose underlying module is locally free of finite type, and
let \(\mathfrak d\) be a Lie subalgebra of \(\mathfrak g\). We call
\(\mathfrak d\) a \emph{Cartan subalgebra} if it satisfies the
following conditions:
\begin{enumerate}
\item[\textup{\SGAInlineRefTarget{sga3:xiv:2.4:item:i}(i)}] \(\mathfrak d\) is locally a direct-factor
  submodule \textup{(and hence is itself locally free of finite
  type)};
\item[\textup{(ii)}] for every \(s\in S\), the Lie algebra
  \(\mathfrak d(s)\) is a Cartan subalgebra of
  \(\mathfrak g(s)\).
\end{enumerate}

\paragraph{\textbf{Definition 2.5.}}
\SGARefTarget{sga3:xiv:2.5}
Let \(S\) and \(\mathfrak g\) be as in \SGARefLink{sga3:xiv:2.4}{2.4}. A section \(a\) of
\(\mathfrak g\) is called \emph{quasi-regular} if, for every
\(s\in S\), the element \(a(s)\in\mathfrak g(s)\) is regular in the
Lie algebra \(\mathfrak g(s)\) over \(k(s)\). We call \(a\) a
\emph{regular section} if it is quasi-regular and is contained in a
Cartan subalgebra of \(\mathfrak g\).

The notions in \SGARefLink{sga3:xiv:2.4}{2.4} and \SGARefLink{sga3:xiv:2.5}{2.5} are again stable under base change and
local for the fpqc topology. Only the last assertion, and only for
the notion of a regular section, requires proof. It follows from the
fact that the Cartan subalgebra containing a given regular section
is uniquely determined. More precisely:
